\section{实际应用：Stable Diffusion 3 与 Flux}

本节讲者简要介绍了 Flow Matching 技术在当前最先进的图像和视频生成模型中的应用。

\subsection{SD3 与 Flux 的流模型基础}

讲者明确指出，许多当前最先进的生成模型已经从扩散模型\textbf{转向}了流模型：

\begin{itemize}
    \item \textbf{Stable Diffusion 3}（SD3）：虽然名字中有``diffusion''，但实际上使用的是 Flow Matching 训练，网络预测的是速度（velocity）而非噪声
    \item \textbf{Flux}：Black Forest Labs 开发的模型，同样基于 Rectified Flow 技术
\end{itemize}

\begin{importantbox}{SD3 的关键设计选择}
Stable Diffusion 3 的核心技术决策包括：
\begin{itemize}
    \item \textbf{Rectified Flow 训练}：使用线性插值条件流映射，网络输出为速度 $v_\theta(\mathbf{x}_t, t)$
    \item \textbf{MM-DiT 架构}：采用 Multimodal Diffusion Transformer 作为骨干网络
    \item \textbf{直线轨迹目标}：通过 Reflow 等技术追求更直的生成轨迹，减少采样步数
    \item \textbf{速度预测（v-prediction）}：而非噪声预测（$\epsilon$-prediction），在某些噪声水平下提供更好的数值稳定性
\end{itemize}
\end{importantbox}

\subsection{Flux 模型}

讲者提到，在课程第一讲中就让学生使用 Flux 模型生成图像。Flux 的一个关键变体是 \textbf{Flux.1-schnell}，它使用了 Rectified Flow 加 Reflow 的技术来实现\textbf{单步}或\textbf{少步}高质量图像生成。

\begin{knowledgebox}{Consistency Model vs.\ Rectified Flow：加速生成的两条技术路线}
在加速扩散/流模型的采样方面，目前存在两条主要的竞争路线：
\begin{itemize}
    \item \textbf{Consistency Model（一致性模型）}：通过 Consistency Distillation (CD) 或 Consistency Training (CT) 学习 ODE 轨迹上的``快捷映射''，直接从任意中间点跳到终点。Lecture 8/9 中讨论过
    \item \textbf{Rectified Flow + Reflow}：通过使轨迹本身变直来减少所需步数。本讲的主题
\end{itemize}

实验比较显示：
\begin{itemize}
    \item Consistency Model 在 1--2 步生成时，ImageNet 64$\times$64 上 FID 约 3--4
    \item 2-Rectified Flow 在 1--2 步生成时，也能达到类似的 FID
\end{itemize}
两种方法各有优劣，目前仍是活跃的研究方向。
\end{knowledgebox}

\subsection{本章小结}

Flow Matching，特别是 Rectified Flow 变体，已经成为当前最先进图像生成模型（SD3、Flux）的基础技术。其核心优势在于训练简洁、轨迹可控、少步生成成为可能。

